\section{Foreign aid: majority support for increased foreign aid if it is well spent}

% AF: ce paragraphe est inutile
Foreign aid is among the most visible and institutionalized forms of international redistribution through which donor countries seek to alleviate poverty and support development abroad. It includes Official Development Assistance (ODA) and related forms of public development finance. Its scale and allocation are largely determined through national legislative and budgetary processes. Sustained increases in foreign aid therefore require sufficient public and legislative support, making donor-country citizens' preferences politically consequential.


Universalism may provide a moral foundation for international redistribution, but policy support can remain limited when citizens are poorly informed or perceive aid as competing with domestic expenditure. % AF: encore une phrase qui n'est pas justifiée par des chiffres ou une citation, à éviter
Information about global inequality may strengthen aid support, while perceived budgetary trade-offs with domestic redistribution may weaken it. International redistribution also presents a collective-action problem: opportunities for countries to free-ride may reduce domestic support when contributions are not internationally coordinated \citep{olson_1965}. % AF: je connais pas ce papier et la phrase ne me permet pas de comprendre ce qu'il faut: est-ce que c'est un papier de théorie qui donne des mécanismes possibles dans l'abstrat ? Quand on cite un papier, il vaut mieux s'en tenir à son contenu: méthodo, résultats.
The literature accordingly identifies moral concern, information, perceived provision, institutional effectiveness, and burden-sharing as distinct determinants of support for cross-border transfers \citep{cappelen_needs_2013,nair_misperceptions_2018,kaufmann_foreign_2019,gampfer_obtaining_2014}. % AF: avant de donner les déterminants du soutien, il faut commencer par donner l'ampleur du soutien (cf. le plan proposé)

The first empirical question is whether national borders constrain distributive concern when potential beneficiaries live abroad and face greater material need. In a real-effort dictator experiment spanning Norway, Germany, Uganda, and Tanzania, \citet{cappelen_needs_2013} find that participants respond to need and entitlement % AF! c'est quoi entitlement? pas clair. Aussi, papier déjà cité au-dessus, faut supprimer l'autre citation
but assign no independent moral weight to nationality. Holding productivity approximately constant, participants from high-income countries allocated 42.7\% to low-income recipients, compared with 36.7\% to high-income recipients. These allocation choices do not, however, establish that citizens support increasing tax-financed foreign-aid budgets; the translation of cross-border concern into policy support depends partly on information. % AF: les deux parties de la phrase ne sont pas logiquement reliées: c'est par parce que l'information influe sur le soutien que celui-ci est majoritaire ou minoritaire. Quel soutien à tax-financed foreign-aid budgets trouvent-ils? Quelle information apportent-ils et qu'est-ce que ça change au soutien ?
In the United States, respondents underestimate their global income rank by 27 percentage points and overestimate median global income by a factor of ten \citep{nair_misperceptions_2018}. % AF: pourquoi le citer plus haut? cette citation me semble mieux
Providing corrective information raises support for increasing foreign aid from 12\% to 22\% and increases donations to foreign charities. % AF: Citer Nair après Kaufmann, dire que l'états-unien médian souhaite maintenir l'aide à son niveau actuel, citer la proportion qui souhaite que l'aide baisse (en plus de hausse), expliquer que cette préférence pour une baisse chez Nair s'explique par une particularité états-unienne et par le fait qu'une majorité soutient une hausse de l'aide conditionnelle
Across nearly 90,000 respondents in 24 donor countries, perceived aid averages 2.65\% of GNI despite actual provision of only 0.37\%; respondents nevertheless prefer aid above its actual level \citep{kaufmann_foreign_2019}. % AF: you mean above its perceived lebel? Pcq y a rien d'étonnant à ce qu'ils préfèrent above actual (je comprends pas le "nevertheless")
% Ce papier il a plein de résultats donc on doit le citer plus en détail, pourquoi pas reproduire leur table principale: 
% Very interesting and useful (above all Table 1 in Appendix). Shows the level of perceived and desired aid in 26 countries between 2005 and 2008. 

% Provide evidence of lobbying by richest against foreign aid (p 19). Indeed, more equal country provide more aid, but because the discrepancy with desired aid is lower, not because people are more generous. This is because of fewer lobbying by richest. (p 4)

% Shows that richer want less aid (p 3), but more unequal countries want more.

% Perceived aid explains 15% of variance of desired aid. Error (desired/perceived) is negatively correlated with income.

% The response is interpreted as ODA although the formulations are (no randomization):  
% (i) What share of national income does your country actually give in foreign aid to help development / poverty alleviation in other countries?  
% (ii) What percentage in your view should it give?

% 11 categories ranging in increasing intervals

% from 0 to higher than 25%: threshold at 0.05; 0.15; 0.35; 0.75; 1.5; 2.5; 4; 7.5; 17.5; 25.

% In most countries (incl. UK, DE, FR, ES but not U.S.) desired aid is larger than perceived. In most countries the gap between the two is small, except in the U.S. where perceived is 7.5% of GDP and preferred is 3%.Use WVS and Gallup (like Chong & Gradstein, Paxton & Knack) but have more waves and the others don't use the question on perceived aid. Shows that richer want less aid ("those in the top income quintile favour ODA (as a share of GNI) that is 0.13 percentage points lower than the preferred share for individuals in the bottom 40\% of the income distribution" after controling for perceived aid).

% Argue that this is due to political influence efforts/possibilities of the rich, as they prefer less aid (support this by a theoretical model + correlations between level of lobbying and actual aid level, controling for desired aid).

Information alone does not determine aid support; institutional design also matters. Evidence from North--South climate finance illustrates these effects particularly clearly. \citet{gampfer_obtaining_2014} find that joint allocation and broader donor burden-sharing increase support in both Germany and the United States, whereas recipient-government efficiency affects support only among US respondents. % AF: pas clair, expliquer plus clairement ce que fait le papier, e.g. Using a conjoint analysis in U.S. and Germany, show xx p.p. stronger support for mitigation/adaptation funding in LIC when cost is shared with other countries (and hence lower), also find preference for recipient countries with efficient administration, and multilateral decision-making => mais il s'agit pas d'aide au dév, si ? changer de section a priori
Compared with recipient-only decision-making, joint allocation raises proposal-selection probabilities by approximately 6 percentage points in the United States and 9 percentage points in Germany. Among US respondents, increasing other donors' contribution from 10\% to 90\% raises the probability that a proposal is selected by 27.9 percentage points. 
Consistent with these findings, 62\% % AF: of whom? say which countries
favor increasing foreign aid, but most condition that increase on aid reaching intended beneficiaries, avoiding diversion, and eliciting greater contributions from other high-income countries \citep{fabre_majority_2025}. % AF: give more detail on these results

Domestic fiscal competition imposes an additional constraint on this conditional support. Economically vulnerable left-leaning respondents remain broadly committed to aid as a moral obligation but are less supportive of increasing public expenditure, a pattern consistent with domestic budgetary trade-offs \citep{cho_ideology_2024}. % AF: rather cite by saying: Eurobarometer shows strong support that foreign aid is a moral obligation and majority support to comply with our promise to increase aid (also questions on donation).
Higher-income respondents prefer lower ODA, while public aid is associated with private giving and actual ODA is negatively correlated with lobbying and private political influence \citep{kaufmann_foreign_2019}. The resulting gap between preferred and actual aid constitutes a democratic aid deficit: public support exists, but it is conditional, fiscally contested, and filtered through unequal political influence.

Foreign aid illustrates the chapter's broader argument: citizens in high-income countries may accept international redistribution, but support depends on credible institutional safeguards and remains vulnerable to domestic fiscal competition. The analysis therefore turns to whether more rule-bound, tax-based instruments can attract equally broad or stronger support.



